\section{Introduction}
\label{sec:introduction}

Distributed systems and microservice architectures rely on Interface Definition Languages such as Google Protocol Buffers \cite{protobuf_spec}, FlatBuffers \cite{flatbuffers} and Cap'n Proto \cite{capnproto}. One schema definition serves several languages, the encoding it produces is compact on the wire, and a message can gain fields without breaking readers built against an earlier version of the schema.

Protocol Buffers, however, depend on an external code generator. Building a C++ project that uses Protobuf requires the build system to invoke \texttt{protoc} and produce intermediate \texttt{.pb.h} and \texttt{.pb.cc} files. This imposes four costs: custom build rules and dependency graphs to track generated artifacts; version skew between the runtime \texttt{libprotobuf} and the build-time \texttt{protoc}, a frequent source of ABI mismatches; codebase bloat, since large systems generate hundreds of thousands of lines that slow incremental builds; and the loss of direct language integration, as schemas remain external build targets rather than first-class compile-time entities.

C++26 makes a different arrangement possible in compilers that implement it. P2996 Static Reflection \cite{p2996} provides programmatic type synthesis through splicing, and P1967 \texttt{\#embed} \cite{p1967} is the standardized way to bring an external file into a translation unit without generating C++ source with tools such as \texttt{xxd -i}. \texttt{ctpb} is built on the former; the evaluated implementation receives the schema text directly rather than through \texttt{\#embed}.

This paper presents \texttt{ctpb} (Compile-Time Protobuf) \cite{ctpb2026}, a header-only C++26 library that carries \texttt{.proto} schemas in the translation unit as compile-time strings, runs a \texttt{constexpr} recursive-descent parser in the compiler front-end, and synthesizes type-safe structs and wire-format routines through static reflection. The contribution is a three-stage architecture that turns raw IDL text into first-class C++ types without generated intermediate files. Its parser handles nested messages, enums, maps, oneofs, and reserved ranges, while P2996 injects fields and makes wire visitors available programmatically. The evaluation reports a $7.487\,\mu\text{s}$ parse for a synthetic 100-field schema and a $4.80\times$ encoding speedup over \texttt{libprotobuf} 3.21.12 on a 47-byte \texttt{Person} workload, with bidirectional wire-format tests against the reference implementation, and it measures the build cost the removed generator is replaced by.
